\documentclass[11pt,border=4pt]{standalone}
% Estilo comum dos esquemáticos do relatório (circuitikz).
% Uso: \documentclass[11pt,border=4pt]{standalone} e, em seguida, % Estilo comum dos esquemáticos do relatório (circuitikz).
% Uso: \documentclass[11pt,border=4pt]{standalone} e, em seguida, % Estilo comum dos esquemáticos do relatório (circuitikz).
% Uso: \documentclass[11pt,border=4pt]{standalone} e, em seguida, \input{estilo_circuitos}.
\usepackage[T1]{fontenc}
\usepackage[utf8]{inputenc}
\usepackage[brazilian]{babel}
\usepackage{lmodern}
\usepackage{amsmath}
\usepackage[american]{circuitikz}
\usetikzlibrary{calc,arrows.meta}

% Cores das figuras originais
\definecolor{anot}{RGB}{0,135,200}      % anotações: valores calculados, setas, resistências vistas
\definecolor{sombra}{RGB}{220,242,253}  % região do transistor nos modelos T e pi
\definecolor{term}{RGB}{62,124,148}     % nomes de terminais destacados (B, C, E)

% Fios com 0,4 pt e corpo dos componentes com 0,8 pt, como nas figuras originais
\ctikzset{bipoles/thickness=2}
.
\usepackage[T1]{fontenc}
\usepackage[utf8]{inputenc}
\usepackage[brazilian]{babel}
\usepackage{lmodern}
\usepackage{amsmath}
\usepackage[american]{circuitikz}
\usetikzlibrary{calc,arrows.meta}

% Cores das figuras originais
\definecolor{anot}{RGB}{0,135,200}      % anotações: valores calculados, setas, resistências vistas
\definecolor{sombra}{RGB}{220,242,253}  % região do transistor nos modelos T e pi
\definecolor{term}{RGB}{62,124,148}     % nomes de terminais destacados (B, C, E)

% Fios com 0,4 pt e corpo dos componentes com 0,8 pt, como nas figuras originais
\ctikzset{bipoles/thickness=2}
.
\usepackage[T1]{fontenc}
\usepackage[utf8]{inputenc}
\usepackage[brazilian]{babel}
\usepackage{lmodern}
\usepackage{amsmath}
\usepackage[american]{circuitikz}
\usetikzlibrary{calc,arrows.meta}

% Cores das figuras originais
\definecolor{anot}{RGB}{0,135,200}      % anotações: valores calculados, setas, resistências vistas
\definecolor{sombra}{RGB}{220,242,253}  % região do transistor nos modelos T e pi
\definecolor{term}{RGB}{62,124,148}     % nomes de terminais destacados (B, C, E)

% Fios com 0,4 pt e corpo dos componentes com 0,8 pt, como nas figuras originais
\ctikzset{bipoles/thickness=2}

% Par diferencial da Questão 3 com entrada diferencial e as fontes c.c. em terra (fig:q3-dif).
% Redesenho de figuras/fig_q3_dif.pdf com folga nos rótulos; entra no relatório em escala 1.
% Compilar com pdflatex.

% Fração das tensões de saída: numerador e denominador um pouco mais afastados da barra,
% para a barra não encostar no índice do numerador nem no denominador.
\newcommand{\fracsaida}[2]{\dfrac{\raisebox{2.3pt}{$#1$}}{\raisebox{-1.1pt}{$#2$}}}

\begin{document}
\begin{circuitikz}
  % Geometria (cm). Nó das fontes em (0,0); colunas dreno-fonte de M1 e M2 em x = -xcol e x = +xcol.
  \def\xcol{2.3}      % meia distância entre as colunas
  \def\ydreno{2.3}    % terminal de dreno dos transistores (centro em ydreno-0,77; fonte em ydreno-1,54)
  \def\ysaida{3.15}   % nó de dreno (saídas)
  \def\ytopo{4.75}    % topo de R_D, ligado à alimentação c.c. em terra
  \def\lsaida{1.0}    % comprimento do fio de saída
  \def\lentrada{0.7}  % comprimento do fio de entrada

  % ---------- transistores (M2 espelhado: porta à direita) ----------
  \node[nmos, arrowmos, anchor=D] (M1) at (-\xcol,\ydreno) {};
  \node[nmos, arrowmos, xscale=-1, anchor=D] (M2) at (\xcol,\ydreno) {};

  % ---------- ramos de dreno: R_D até a alimentação c.c. em terra ----------
  \draw (M1.D) -- (-\xcol,\ysaida) to[R] (-\xcol,\ytopo) node[ground, rotate=180] {};
  \draw (M2.D) -- (\xcol,\ysaida) to[R] (\xcol,\ytopo) node[ground, rotate=180] {};

  % ---------- fontes interligadas e R_S até a terra ----------
  \draw (M1.S) -- (-\xcol,0) -- (\xcol,0) -- (M2.S);
  \draw (0,0) -- (0,-0.45) to[R] (0,-2.05) node[ground] {};
  \node[circ] at (0,0) {};

  % ---------- saídas ----------
  \draw (-\xcol,\ysaida) -- (-\xcol-\lsaida,\ysaida) node[ocirc] {};
  \draw (\xcol,\ysaida) -- (\xcol+\lsaida,\ysaida) node[ocirc] {};
  \node[circ] at (-\xcol,\ysaida) {};
  \node[circ] at (\xcol,\ysaida) {};

  % ---------- entradas ----------
  \draw (M1.G) -- ++(-\lentrada,0) node[ocirc] (E1) {};
  \draw (M2.G) -- ++(\lentrada,0) node[ocirc] (E2) {};

  % ---------- nomes dos componentes ----------
  \node[anchor=east, inner sep=0pt] at (-\xcol-0.36,{(\ysaida+\ytopo)/2}) {$R_D$};
  \node[anchor=west, inner sep=0pt] at (\xcol+0.36,{(\ysaida+\ytopo)/2}) {$R_D$};
  \node[anchor=east, inner sep=0pt] at (-0.36,-1.25) {$R_S$};
  \node[anchor=west, inner sep=0pt] at (-\xcol+0.3,\ydreno-0.77) {$M_1$};
  \node[anchor=east, inner sep=0pt] at (\xcol-0.3,\ydreno-0.77) {$M_2$};

  % ---------- anotações ----------
  \begin{scope}[anot]
    % tensões de saída
    \node[anchor=east, inner sep=0pt] at (-\xcol-\lsaida-0.2,\ysaida)
      {$v_{o1} = -g_m(R_D \parallel r_o)\:\fracsaida{v_{id}}{2}$};
    \node[anchor=west, inner sep=0pt] at (\xcol+\lsaida+0.2,\ysaida)
      {$v_{o2} = +g_m(R_D \parallel r_o)\:\fracsaida{v_{id}}{2}$};
    % correntes de sinal nos drenos (para baixo em M1, para cima em M2)
    \draw[line width=0.8pt, -latex] (-\xcol+0.3,{(\ydreno+\ysaida)/2+0.25}) -- (-\xcol+0.3,{(\ydreno+\ysaida)/2-0.25});
    \node[anchor=west, inner sep=0pt] at (-\xcol+0.52,{(\ydreno+\ysaida)/2}) {$g_m v_{id}/2$};
    \draw[line width=0.8pt, -latex] (\xcol-0.3,{(\ydreno+\ysaida)/2-0.25}) -- (\xcol-0.3,{(\ydreno+\ysaida)/2+0.25});
    \node[anchor=east, inner sep=0pt] at (\xcol-0.52,{(\ydreno+\ysaida)/2}) {$g_m v_{id}/2$};
    % excitação diferencial
    \node[anchor=east, inner sep=0pt] at ($(E1)+(-0.2,0)$) {$+v_{id}/2$};
    \node[anchor=west, inner sep=0pt] at ($(E2)+(0.2,0)$) {$-v_{id}/2$};
    % tensões porta-fonte, sob os terminais de entrada
    \node[align=center, inner sep=0pt] at (-3.9,0.3) {$+$\\[-2pt] $v_{gs1} = v_{id}/2$\\[-2pt] $-$};
    \node[align=center, inner sep=0pt] at (3.9,0.3) {$+$\\[-2pt] $v_{gs2} = -v_{id}/2$\\[-2pt] $-$};
    % nó das fontes e R_S
    \node[anchor=north west, inner sep=0pt] at (0.2,-0.2) {$0\,\mathrm{V}$ (terra virtual)};
    \node[anchor=west, inner sep=0pt] at (0.42,-1.25) {sem corrente de sinal};
  \end{scope}
\end{circuitikz}
\end{document}
